\section{Introduction}

Large Language Models (LLMs) have revolutionized the field of AI technologies, providing powerful tools for automated reasoning and conversational interactions~\cite{yang2025qwen3technicalreport,deepseekai2025deepseekv3technicalreport,5team2025glm45agenticreasoningcoding}. Their strengths lie in generative fluency and contextual understanding. However, these models face fundamental challenges when dealing with fact-rich domains there knowledge consistency, scalability and groundendness are crucial. Integration of external knowledge graphs (KGs) into LLM-driven systems gives a promising opportunity to bridge the gaps between reasoning and factuality~\cite{chepurova-etal-2026-wikontic,bai2025autoschemakgautonomousknowledgegraph,10.1145/3746027.3755628}. Yet, the complexity of scaling KG-based methods for open-domain QA tasks and maintaining high retrieval precision remains a bottleneck.

Graph-based Retrieval-Augmented Generation (GraphRAG)~\cite{Gao2023RetrievalAugmentedGF , 10.1145/3777378} frameworks have gained prominence by augmenting prompts with retrieved information, yet they remain restricted by static ontology and inefficient traversal mechanisms. Thus, dynamic, tailored algorithms for knowledge retrieval and reasoning are crucial for maximizing the utility of KGs in combination with LLMs. Traditional GraphRAG systems rely predominantly on node-level retrievals, limiting their scalability and precision~\cite{hu-etal-2025-grag,mavromatis-karypis-2025-gnn,luo2024graph}. They face difficulties in handling multi-hop reasoning tasks, where search strategy must be dynamic and modify based on intermediate discovered information. Further, static retrieval patterns limit their adaptability to varied domains and user intents.

To address these challenges, we propose PersonalAI 2.0 (PAI-2), a GraphRAG method that incorporates graph-based external memory to store unstructured textual knowledge alongside LM-driven reasoning. By introducing a multi-stage query-processing pipeline, PAI-2 aims to optimize graph traversal and query resolution. Its contributions lie in dynamically planned, iterative information searches, guided by entity extraction and vertex matching. By systematically decomposing complex queries into manageable subqueries, PAI-2 ensures focused retrieval of only relevant segments of underlying knowledge graph. Ultimately, this modification holds promise for improving factuality and reducing hallucinations across multi-hop reasoning tasks.

Proposed method can be applied in a wide range of fields: from personalized education platforms to customer service chatbots, where contextual awareness and precision are highly important. Beyond theoretical advancement, PAI-2 lays foundational principles for designing future-generation LLMs, augmented with richer, structured external memory graphs.

In summary, our main contributions are as follows:
\begin{enumerate}
    \item We propose PersonalAI 2.0 (PAI-2), a GraphRAG method which effectively integrates graph based external memory to store unstructured knowledge from texts and LLM reasoning abilities to plan information search and manage/specify graph traversal.
    \item We evaluate PAI-2 on Natural Questions, TriviaQA, HotpotQA, 2WikiMultihopQA, MuSiQue, DiaASQ benchmarks and compare it with LightRAG, RAPTOR, HippoRAG 2. Our method shows superior performance on 4 out of 6 benchmarks with average gain 4\% by LLM-as-a-Judge. 
    \item We show that plan enhancing mechanism during information search increases answer accuracy on average 18\% by LLM-as-a-Judge across six datasets.
    \item We show that use of graph traversal algorithms (e.g. Beam Search, WaterCircles) gains superior performance compared to standard flatten retriever: on average 6\% by LLM-as-a-Judge across six datasets.
    \item PAI-2 achieves state-of-the-art results on the MINE-1 benchmark, reaching 89\% information-retention score. We show that PAI`s memory construction algorithm is more stable (less LLM parsing errors), compared to KGGen and Wikontic in 7-14B LLM setting. 
\end{enumerate}
%- поиск информации с планированием
%- SOTA среди отобранных аналогов
%- показано, что корректировка плана поиска даёт прирост эффективности
%- показано, что использование алгоритмов обхода графа даёт прирост эффективности по сравнения с плским поиском 
%- более стабильный механизм формирования пространства памяти среди аналогов
% TODO